\problemname{Dalarna} % storyn skulle absolut kunna bytas ut

\noindent
Harry älskar Dalarna och drömmer om att det ska bli Sveriges högst rankade län. Efter noggrann analys har Harry 
kommit fram till att om varje stad $i$ i Dalarna hade $p_i$ fler bilresor genom staden skulle Dalarna ta 
förstaplatsen. För att uppnå detta planerar Harry att beställa taxibilar som kör mellan städerna där han kan 
beställa att en bil ska ta kortaste vägen mellan två städer.

Dalarnas vägnät kan representeras som ett träd med $N$ städer. För varje stad $i$ behöver Harry att minst 
$p_i$ taxibilar passerar genom staden för att nå sitt mål. En taxibil kan köra längs en godtycklig väg i trädet.

För att inte väcka misstankar vill Harry använda så få taxibilar som möjligt. Hjälp Harry att räkna ut det 
minsta antalet taxibilar han behöver beställa för att uppnå sitt mål.

\section*{Indata}
Den första raden innehåller ett heltal $N$ ($1 \le N \le 5 \cdot 10^5$), antalet städer i Dalarna.

Den andra raden innehåller $N$ heltal $p_1, p_2, ..., p_N$ ($0 \le p_i \le 10^9$), där $p_i$
anger hur många taxibilar som måste passera genom stad $i$.

De följande $N-1$ raderna beskriver vägarna mellan städerna. Varje rad innehåller två heltal $u$ och $v$
($1 \le u, v \le N$), vilket betyder att det finns en väg mellan stad $u$ och stad $v$.
Vägarna bildar ett träd.

\section*{Utdata}
Skriv ut ett heltal: det minsta antalet taxibilar som Harry behöver beställa.

\section*{Poängsättning}
Din lösning kommer att testas på en mängd testfallsgrupper.
För att få poäng för en grupp så måste du klara alla testfall i gruppen.

\noindent
\begin{tabular}{| l | l | p{12cm} |}
  \hline
  \textbf{Grupp} & \textbf{Poäng} & \textbf{Gränser} \\ \hline
  $1$    & $8$        & Trädet är en linje. \\ \hline
  $2$    & $11$       & $p_i = 1$ för alla $i$. \\ \hline
  $3$    & $9$        & $p_i \le 1$ för alla $i$. \\ \hline
  $4$    & $12$       & $N \le 7$, $p_i \le 5$ för alla $i$. \\ \hline
  $5$    & $14$       & $N \le 100$, $p_i \le 100$ för alla $i$. \\ \hline
  $6$    & $15$       & $N \le 1000$ \\ \hline
  $7$    & $31$       & Inga ytterligare begränsningar. \\ \hline
\end{tabular}
